On considère la fonction $f$ définie sur l'intervalle $]0~;~ + \infty[$ par

\[f(x) = \frac{\e^{\sqrt{x}}}{2\sqrt{x}}\]

et on appelle $\mathcal{C}_f$ sa courbe représentative dans un repère orthonormé.

\medskip

\begin{enumerate}
\item On définit la fonction $g$ sur l'intervalle $]0~;~ + \infty[$ par $g(x) = \e^{\sqrt{x}}$.
	\begin{enumerate}
		\item Montrer que $g'(x) = f(x)$ pour tout $x$ de l'intervalle $]0~;~ + \infty[$.
		\item Pour tout réel $x$ de l'intervalle $]0~;~ + \infty[$, calculer $f'(x)$ et montrer que:

\[f'(x) = \dfrac{\e^{\sqrt{x}}\left(\sqrt x - 1\right)}{4x\sqrt x}.\]
	\end{enumerate}
\item
	\begin{enumerate}
		\item Déterminer la limite de la fonction $f$ en $0$.
		\item Interpréter graphiquement ce résultat.
	\end{enumerate}
\item
	\begin{enumerate}
		\item Déterminer la limite de la fonction $f$ en $+ \infty$.
		\item Étudier le sens de variation de la fonction $f$ sur $]0~;~+ \infty[$.

Dresser le tableau de variations de la fonction $f$ en y faisant figurer les limites aux bornes de l'intervalle de définition.
		\item Montrer que l'équation $f(x) = 2$ admet une unique solution sur l'intervalle $[1~;~+\infty[$ et donner une valeur approchée à $10^{-1}$ près de cette solution.
	\end{enumerate}
\item On pose  $I = \displaystyle\int_1^2 f(x)\, \text{d}x$.
	\begin{enumerate}
		\item Calculer $I$.
		\item Interpréter graphiquement le résultat.
	\end{enumerate}
\item On admet que la fonction $f$ est deux fois dérivable sur l'intervalle $]0~;~ + \infty[$ et que:

\[f"(x) = \dfrac{\e^{\sqrt{x}}\left(x - 3\sqrt x + 3\right)}{8x^2\sqrt x}.\]

	\begin{enumerate}
		\item En posant $X = \sqrt x$, montrer que $x - 3\sqrt x + 3 > 0$ pour tout réel $x$ de l'intervalle $]0~;~ + \infty[$.
		\item Étudier la convexité de la fonction $f$ sur l'intervalle $]0~;~ + \infty[$.
	\end{enumerate}
\end{enumerate}
